% !TEX root = master.tex
% !TEX program = lualatex

\chapter{Basics - 2}%
\label{cha:Basics - 2}

\section{C01b — Basics 6 : Tableaux}

\subsection{Attention générale}

Les tableaux en C ressemblent syntaxiquement à ceux de Java,  
mais leur comportement est très différent :

\begin{itemize}
    \item ce ne sont pas des objets
    \item il n’y a aucune vérification automatique
    \item il n’y a pas d’exception en cas de débordement
    \item la taille est généralement fixe
\end{itemize}

C’est une source majeure d’erreurs pour les programmeurs venant de Java.

\subsection{Taille des tableaux}

En C classique :

\begin{itemize}
    \item la taille doit être connue à la compilation
    \item on ne peut pas redimensionner un tableau
\end{itemize}

Pour des tailles dynamiques, on utilisera les pointeurs
(allocation dynamique, vu plus tard).

\subsection{Déclaration}

Syntaxe identique à Java :

\begin{lstlisting}[language=C]
int t[5];
\end{lstlisting}

Mais il est recommandé d’éviter les nombres magiques :

\begin{lstlisting}[language=C]
#define N 5
int t[N];
\end{lstlisting}

Ou avec une constante :

\begin{lstlisting}[language=C]
const int N = 5;
int t[N];
\end{lstlisting}

\subsection{Variable Length Arrays (VLA)}

Depuis C99, on peut déclarer un tableau dont la taille est connue
à l’exécution :

\begin{lstlisting}[language=C]
int n;
scanf("%d",&n);
int t[n];
\end{lstlisting}

La taille n’est pas connue à la compilation,
mais une fois fixée, elle ne peut plus changer.

\subsection{Initialisation}

Comme en Java :

\begin{lstlisting}[language=C]
int t[5] = {1,2,3,4,5};
\end{lstlisting}

On peut aussi laisser le compilateur déduire la taille :

\begin{lstlisting}[language=C]
int t[] = {1,2,3,4,5};
\end{lstlisting}

Cela signifie simplement que la taille est le nombre d’éléments.

\subsection{Initialisation à zéro}

Si au moins une valeur est fournie,
toutes les autres sont initialisées à zéro :

\begin{lstlisting}[language=C]
int t[1000] = {0};   // tout le tableau vaut 0
\end{lstlisting}

\subsection{Initialisation partielle (C99)}

On peut initialiser une position précise :

\begin{lstlisting}[language=C]
int t[10] = {[2] = 5};
\end{lstlisting}

Résultat :

\[
t = (0,0,5,0,0,0,0,0,0,0)
\]

On peut combiner plusieurs initialisations :

\begin{lstlisting}[language=C]
int t[10] = {1,2,3,[7]=9};
\end{lstlisting}

Les valeurs non spécifiées valent 0.

\subsection{Accès aux éléments}

Syntaxe identique à Java :

\begin{lstlisting}[language=C]
t[i]
\end{lstlisting}

Mais :

\begin{itemize}
    \item aucun contrôle de débordement
    \item accès hors limites = comportement indéfini
\end{itemize}

\textbf{Règle : indices valides de 0 à taille-1.}

\subsection{Un tableau ne connaît jamais sa taille}

Un tableau C :

\begin{itemize}
    \item ne stocke pas sa taille
    \item ne peut pas la fournir
    \item le programmeur doit la connaître
\end{itemize}

\subsection{Passage d’un tableau à une fonction}

Lorsqu’on passe un tableau à une fonction :

\begin{itemize}
    \item il est converti en pointeur
    \item la taille n’est pas transmise
\end{itemize}

Donc :

\begin{lstlisting}[language=C]
void f(int t[], int n);
\end{lstlisting}

équivaut à :

\begin{lstlisting}[language=C]
void f(int *t, int n);
\end{lstlisting}

\subsection{Règle absolue}

\[
\textbf{Toujours passer la taille avec un tableau.}
\]

\subsection{Tableaux passés par référence}

Même sans \texttt{\&}, un tableau est toujours transmis
par référence (via un pointeur).

Exemple :

\begin{lstlisting}[language=C]
void g(int t[], int n) {
    for(int i=0;i<n;i++)
        t[i]=i+1;
}
\end{lstlisting}

Dans \texttt{main} :

\begin{lstlisting}[language=C]
int t[3]={0,0,0};
g(t,3);
\end{lstlisting}

Après l’appel :

\[
t = (1,2,3)
\]

Donc la fonction modifie réellement le tableau.

\subsection{Empêcher la modification}

On peut protéger le tableau avec \texttt{const} :

\begin{lstlisting}[language=C]
void g(const int t[], int n);
\end{lstlisting}

Toute modification produira une erreur de compilation.

\subsection{Retourner un tableau}

Impossible de retourner un tableau directement :

\begin{itemize}
    \item pas d’affectation globale de tableau
    \item pas de retour par valeur
\end{itemize}

Il faudra utiliser :

\begin{itemize}
    \item un pointeur
    \item ou une structure
\end{itemize}

\subsection{Tableaux multidimensionnels}

Même principe qu’en Java :

\begin{lstlisting}[language=C]
int m[3][4];
\end{lstlisting}

Mais lorsqu’on passe un tableau à une fonction,
il faut spécifier les dimensions intermédiaires :

\begin{lstlisting}[language=C]
void f(int m[][4], int n);
\end{lstlisting}

Cela permet au compilateur de calculer les adresses mémoire.

Cependant :

\begin{itemize}
    \item le tableau ne connaît toujours pas sa taille
    \item il faut passer explicitement lignes et colonnes
\end{itemize}

\section{C01b — Basics 7 : enum, typedef et struct}

\subsection{Types énumérés (enum)}

Un type énuméré permet de définir un nouveau type constitué d’un ensemble
de valeurs nommées.

Contrairement à Java :
\begin{itemize}
    \item ce n’est pas un objet
    \item c’est essentiellement un entier nommé
\end{itemize}

\subsubsection{Définition}

\begin{lstlisting}[language=C]
enum couleur {BLANC, NOIR, VIDE};
\end{lstlisting}

Déclaration :

\begin{lstlisting}[language=C]
enum couleur c = BLANC;
\end{lstlisting}

Le vrai nom du type est :
\[
\texttt{enum couleur}
\]

\subsubsection{Conversion avec int}

En C, la conversion entre enum et int est implicite :

\begin{itemize}
    \item enum → int autorisé
    \item int → enum autorisé
\end{itemize}

Cela rend les enums plus flexibles mais moins sûrs qu’en Java.

\subsubsection{Pattern utile : nombre d’éléments}

On peut ajouter une dernière valeur :

\begin{lstlisting}[language=C]
enum canton {
    VAUD, VALAIS, GENEVE, JURA,
    NB_CANTONS
};
\end{lstlisting}

Cela permet d’utiliser \texttt{NB\_CANTONS} comme taille.

\subsection{Alias de type (typedef)}

\texttt{typedef} permet de donner un nouveau nom à un type existant.

\subsubsection{Pourquoi utiliser typedef}

\begin{itemize}
    \item nommer un concept métier
    \item clarifier les signatures de fonctions
    \item faciliter un changement de type ultérieur
\end{itemize}

\subsubsection{Exemple simple}

\begin{lstlisting}[language=C]
typedef int distance;
distance d = 3;
\end{lstlisting}

Si on change plus tard \texttt{int} en \texttt{double},
une seule ligne est modifiée.

\subsubsection{Exemple conceptuel}

\begin{lstlisting}[language=C]
typedef double vecteur[3];
vecteur v;
\end{lstlisting}

Cela permet de dire explicitement que la fonction manipule
un vecteur plutôt qu’un simple tableau.

\subsubsection{Règle pratique}

Pour écrire un typedef :

\begin{itemize}
    \item écrire une déclaration normale
    \item remplacer le nom de variable par le nom du type
    \item ajouter \texttt{typedef} devant
\end{itemize}

\subsection{Structures (struct)}

Une structure permet de regrouper des données hétérogènes.

On peut voir une struct comme :

\begin{itemize}
    \item un objet sans méthodes
    \item dont tous les champs sont publics
\end{itemize}

\subsubsection{Définition}

\begin{lstlisting}[language=C]
struct personne {
    char nom[50];
    int age;
    double taille;
    char sexe;
};
\end{lstlisting}

Déclaration :

\begin{lstlisting}[language=C]
struct personne p;
\end{lstlisting}

Le nom du type est :
\[
\texttt{struct personne}
\]

\subsection{Structs imbriquées}

On peut utiliser une struct dans une autre :

\begin{lstlisting}[language=C]
struct point {
    double x, y;
};

struct cercle {
    struct point centre;
    double rayon;
};
\end{lstlisting}

\subsection{Struct + typedef (pattern courant)}

Pour éviter d’écrire \texttt{struct} à chaque fois :

\begin{lstlisting}[language=C]
typedef struct {
    char nom[50];
    int age;
    double taille;
} Personne;
\end{lstlisting}

On peut alors écrire :

\begin{lstlisting}[language=C]
Personne p;
\end{lstlisting}

\subsection{Initialisation}

Comme pour les tableaux :

\begin{lstlisting}[language=C]
Personne p = {"Jean", 20, 1.75};
\end{lstlisting}

Depuis C99, initialisation partielle :

\begin{lstlisting}[language=C]
Personne p = {.age = 20, .taille = 1.75};
\end{lstlisting}

Bonne pratique :
\begin{lstlisting}[language=C]
Personne p = {0};
\end{lstlisting}
pour initialiser toute la structure à zéro.

\subsection{Accès aux champs}

Syntaxe identique à Java :

\begin{lstlisting}[language=C]
p.age = 20;
printf("%d", p.age);
\end{lstlisting}

\subsection{Passage par pointeur}

Si on passe une structure par pointeur :

\begin{lstlisting}[language=C]
void vieillir(Personne *p) {
    p->age++;
}
\end{lstlisting}

L’opérateur \texttt{->} signifie :

\[
p->age \equiv (*p).age
\]

Car comme on manipule une adresse on ne peut plus utiliser p.age .

\subsection{Retour d’une structure}

Contrairement aux tableaux,
une structure peut être retournée :

\begin{lstlisting}[language=C]
Personne naissance() {
    Personne p = {0};
    return p;
}
\end{lstlisting}

\subsection{Affectation globale}

On peut copier une structure entière :

\begin{lstlisting}[language=C]
p2 = p1;
\end{lstlisting}

Cela copie tous les champs.

\subsection{Comparaison}

Il n’existe pas d’opérateur global :

\begin{lstlisting}[language=C]
p1 == p2   // interdit
\end{lstlisting}

Il faut écrire une fonction de comparaison.

\subsection{Tableaux de structures}

\begin{lstlisting}[language=C]
Personne groupe[10];
\end{lstlisting}

On peut initialiser chaque élément individuellement.

\subsection{Flexible array member (avancé)}

Depuis C99, une struct peut se terminer par un tableau
de taille variable :

\begin{lstlisting}[language=C]
struct data {
    int n;
    int values[];
};
\end{lstlisting}

Ce tableau doit être le dernier champ et sera géré
par allocation dynamique (vu avec les pointeurs).

\section{C01b — Basics 8 : Types et mots-clés avancés}

\subsection{Objectif}

Cette section introduit quelques éléments plus avancés du langage C :

\begin{itemize}
    \item modificateurs supplémentaires
    \item unions
    \item champs de bits
    \item structures/unions anonymes (C11)
\end{itemize}

Ces éléments ne sont pas indispensables pour débuter en C,
mais il est utile de savoir les reconnaître.

\subsection{Rappel des modificateurs déjà vus}

On a déjà rencontré :

\begin{itemize}
    \item \texttt{const}
    \item \texttt{long}
    \item \texttt{short}
    \item \texttt{unsigned}
\end{itemize}

Le modificateur \texttt{signed} existe mais correspond au comportement par défaut,
donc il est rarement utilisé.

\subsection{Le mot-clé \texttt{auto}}

\texttt{auto} indique qu’une variable est locale à un bloc.

Cependant, c’est le comportement par défaut en C moderne,
donc ce mot-clé n’est quasiment jamais utilisé.

\subsection{Le mot-clé \texttt{static}}

Le mot-clé \texttt{static} a deux significations différentes.

\subsubsection{Static au niveau fichier}

Il limite la visibilité d’une variable ou fonction au fichier courant.

Cela sera surtout important lors de la compilation séparée.

\subsubsection{Static dans une fonction}

Une variable \texttt{static} locale :

\begin{itemize}
    \item est initialisée une seule fois
    \item conserve sa valeur entre les appels
\end{itemize}

\begin{lstlisting}[language=C]
void f(void) {
    static int compteur = 0;
    compteur++;
    printf("%d\n", compteur);
}
\end{lstlisting}

Chaque appel incrémente la même variable.

On peut voir cela comme une mémoire interne à la fonction.

\subsection{Le mot-clé \texttt{register}}

\texttt{register} suggère au compilateur de stocker une variable
dans un registre processeur plutôt qu’en mémoire.

Cependant :

\begin{itemize}
    \item ce n’est qu’une suggestion
    \item les compilateurs modernes optimisent déjà automatiquement
\end{itemize}

Donc ce mot-clé est aujourd’hui rarement utile.

\subsection{Le mot-clé \texttt{restrict}}

\texttt{restrict} s’applique aux pointeurs.

Il signifie que ce pointeur est le seul à référencer cette zone mémoire.

Cela permet au compilateur d’effectuer davantage d’optimisations.

Ce mot-clé sera compris après l’étude des pointeurs.

\subsection{Le mot-clé \texttt{volatile}}

\texttt{volatile} indique que la valeur peut changer indépendamment du programme.

Exemples :

\begin{itemize}
    \item registre matériel
    \item capteur externe
    \item mémoire partagée
\end{itemize}

Le compilateur n’effectuera alors aucune optimisation
supposant que la valeur reste constante.

\subsection{Les unions}

Une union ressemble à une structure,
mais ses champs partagent la même zone mémoire.

\subsubsection{Principe}

\begin{lstlisting}[language=C]
union valeur {
    int i;
    double d;
};
\end{lstlisting}

Une variable peut contenir :

\begin{itemize}
    \item soit un entier
    \item soit un double
    \item mais pas les deux en même temps
\end{itemize}

\subsubsection{Problème principal}

Il faut toujours savoir quel champ est valide.

\begin{lstlisting}[language=C]
union valeur x;
x.i = 3;
printf("%f", x.d);   // valeur incohérente
\end{lstlisting}

La mémoire est réinterprétée selon le type utilisé.

\subsubsection{Pourquoi éviter les unions}

\begin{itemize}
    \item difficile à utiliser correctement
    \item facile de produire des incohérences
    \item souvent remplaçables par des pointeurs génériques
\end{itemize}

\subsubsection{Petit avantage}

Une union alloue automatiquement la taille du plus grand champ,
ce qui peut économiser de la mémoire.

\subsection{Structures et unions anonymes (C11)}

Depuis C11, on peut imbriquer des structures ou unions sans nom.

\begin{lstlisting}[language=C]
struct exemple {
    int cle;
    union {
        int i;
        double d;
    };
};
\end{lstlisting}

On peut ensuite écrire directement :

\begin{lstlisting}[language=C]
x.i = 3;
\end{lstlisting}

sans nom intermédiaire.

\subsection{Champs de bits}

Les champs de bits permettent de spécifier le nombre exact de bits utilisés.

\begin{lstlisting}[language=C]
struct nombre {
    unsigned signe : 1;
    unsigned exposant : 15;
    unsigned mantisse : 32;
};
\end{lstlisting}

Cela permet de créer des représentations compactes en mémoire.

Cependant :

\begin{itemize}
    \item c’est très dépendant de l’architecture
    \item c’est rarement utilisé dans les programmes classiques
\end{itemize}

\section{C01b — Basics 9 : Entrées / sorties standard}

\subsection{Principe des entrées/sorties}

Les entrées/sorties permettent au programme de communiquer avec l’extérieur.

En C, elles se font via des \textbf{flots} (\textit{streams}) :
des variables représentant un canal d’échange avec l’extérieur.

Les flots standards sont :

\begin{itemize}
    \item \texttt{stdin} : entrée standard (clavier)
    \item \texttt{stdout} : sortie standard (terminal)
    \item \texttt{stderr} : sortie d’erreur standard
\end{itemize}

Ces flots sont définis dans :

\begin{lstlisting}[language=C]
#include <stdio.h>
\end{lstlisting}

\subsection{La fonction printf}

\texttt{printf} écrit sur \texttt{stdout}.

Prototype :

\begin{lstlisting}[language=C]
int printf(const char *format, ...);
\end{lstlisting}

Elle retourne :

\begin{itemize}
    \item le nombre de caractères écrits
    \item une valeur négative en cas d’erreur
\end{itemize}

\subsection{Chaîne de format}

Une chaîne de format contient :

\begin{itemize}
    \item du texte normal (écrit tel quel)
    \item des \textbf{placeholders} introduits par \%
\end{itemize}

Structure :

\[
\% \; [modificateurs] \; type
\]

Exemples :

\begin{itemize}
    \item \texttt{\%d} : entier
    \item \texttt{\%u} : entier non signé
    \item \texttt{\%f} : double
    \item \texttt{\%s} : chaîne
    \item \texttt{\%p} : pointeur
    \item \texttt{\%\%} : caractère \%
\end{itemize}

\subsection{Largeur et précision}

Exemple :

\begin{lstlisting}[language=C]
printf("%5.2f", x);
\end{lstlisting}

\begin{itemize}
    \item 5 = nombre total de caractères (espaces et point inclus) que l'affichage doit occuper au minimum. Si le nombre est plus petit, le C ajoute des espaces vides à gauche pour "pousser" le chiffre.
    \item .2 = 2 décimales
\end{itemize}

Si la valeur dépasse la largeur, elle n’est pas tronquée.

\subsection{Modificateurs utiles}

\begin{itemize}
    \item \texttt{-} : alignement à gauche
    \item \texttt{+} : afficher le signe
    \item \texttt{0} : remplissage par des zéros
    \item \texttt{\#} : forcer certains formats
\end{itemize}

\subsection{Formats flottants}

\begin{itemize}
    \item \texttt{\%f}(Fixed): Il affiche le nombre tel quel, avec des zéros inutiles à la fin si nécessaire (par défaut 6 chiffres après la virgule).
    \item \texttt{\%g}(General/Adaptatif) : C'est le mode "intelligent". Il choisit le format le plus court. Il supprime les zéros inutiles à la fin et peut passer en notation scientifique (1e+05) si le nombre est trop grand ou trop petit.
\end{itemize}

\subsection{Bufferisation de stdout}

\texttt{stdout} est bufferisé :

\begin{itemize}
    \item \texttt{printf} n’affiche pas immédiatement
    \item il écrit dans un tampon
    \item le système l’affiche plus tard
\end{itemize}

Donc un message peut ne jamais apparaître si le programme
crashe avant le flush.

\subsection{Forcer l’affichage}

Deux méthodes :

\begin{itemize}
    \item ajouter \texttt{\textbackslash n}
    \item utiliser \texttt{fflush(stdout);}
\end{itemize}

\begin{lstlisting}[language=C]
printf("Hello");
fflush(stdout);
\end{lstlisting}

\textbf{Utiliser fflush avec parcimonie.}

\subsection{Lecture avec scanf}

\texttt{scanf} lit sur \texttt{stdin}.

Prototype :

\begin{lstlisting}[language=C]
int scanf(const char *format, ...);
\end{lstlisting}

Elle retourne :

\begin{itemize}
    \item le nombre de champs lus correctement
    \item 0 en cas d’échec
\end{itemize}

\subsection{Important : scanf utilise des pointeurs}

Exemple :

\begin{lstlisting}[language=C]
int x;
scanf("%d", &x);
\end{lstlisting}

On passe l’adresse de la variable.

\subsection{Lecture arrêtée par les blancs}

\texttt{scanf} s’arrête au premier espace, tabulation ou retour ligne.

Donc :

\begin{itemize}
    \item impossible de lire une phrase complète avec \texttt{\%s}
    \item préférable d’utiliser \texttt{fgets}
\end{itemize}

\subsection{Lire une chaîne avec fgets}

\begin{lstlisting}[language=C]
char buf[100];
fgets(buf, 100, stdin);
\end{lstlisting}

\texttt{fgets} :

\begin{itemize}
    \item lit toute la ligne
    \item inclut le \texttt{\textbackslash n}
    \item évite les dépassements mémoire
\end{itemize}

\subsection{Format spécial de scanf}

\textbf{Jeu de caractères}

\begin{lstlisting}[language=C]
scanf("%[A-Z]", buf);
\end{lstlisting}

lit uniquement les majuscules.

\textbf{Négation}

\begin{lstlisting}[language=C]
scanf("%[^\n]", buf);
\end{lstlisting}

lit tout sauf le retour ligne.

\subsection{Lire un double}

Important :

\begin{itemize}
    \item afficher un double : \texttt{printf} : \texttt{\%f}
    \item lire un double : \texttt{scanf} : \texttt{\%lf}
\end{itemize}

\subsection{Ignorer un champ}

\begin{lstlisting}[language=C]
scanf("%d %*d %lf", &i, &x);
\end{lstlisting}

Le champ avec \texttt{*} est lu mais ignoré.

\subsection{Problème classique : boucle infinie}

Si \texttt{scanf} échoue (par exemple si on demande un int et que l'utilisateur donne une chaîne de caractère):

\begin{itemize}
    \item la variable ne change pas
    \item l’entrée reste dans le buffer \textrightarrow on va la relire au prochain scanf
    \item la boucle recommence infiniment 
\end{itemize}

\subsection{Solution robuste}

Toujours tester le retour :

\begin{lstlisting}[language=C]
int r = scanf("%d", &i);
if(r != 1) { //car scanf renvoie le nombre d'éléments qu'il a lu
    printf("Erreur\n");
}
\end{lstlisting}

Et vider le buffer (passer à la prochaine ligne) :

\begin{lstlisting}[language=C]
int c;
// On lit chaque caractère jusqu'à la fin de la ligne \n ou du fichier (EOF)
while((c=getchar())!='\n' && c!=EOF);
\end{lstlisting}

\subsection{stderr}

\texttt{stderr} est un flot séparé :

\begin{itemize}
    \item destiné aux messages d’erreur
    \item non bufferisé
    \item affiché immédiatement
\end{itemize}

Pour écrire dessus :

\begin{lstlisting}[language=C]
fprintf(stderr, "Erreur\n");
\end{lstlisting}

Bonne pratique :

\begin{itemize}
    \item messages normaux → stdout
    \item erreurs → stderr
\end{itemize}

\section{C01b — Basics 10 : Fichiers}

\subsection{Principe}

Un fichier est manipulé en C via un \textbf{flot} (\texttt{FILE}).

On ne manipule jamais un fichier directement :
on manipule un pointeur vers un objet de type \texttt{FILE}.

\begin{lstlisting}[language=C]
#include <stdio.h>

FILE *f;
\end{lstlisting}

Les fichiers sont toujours manipulés par référence (via pointeur).

\subsection{Ouverture d’un fichier}

Fonction principale :

\begin{lstlisting}[language=C]
FILE *fopen(const char *nom, const char *mode);
\end{lstlisting}

Exemple :

\begin{lstlisting}[language=C]
FILE *entree = fopen("data.txt", "r");
\end{lstlisting}

\subsection{Modes principaux}

\begin{itemize}
    \item \texttt{"r"} : lecture
    \item \texttt{"w"} : écriture (écrase le fichier)
    \item \texttt{"a"} : écriture en fin de fichier
\end{itemize}

\subsection{Modificateurs}

\begin{itemize}
    \item \texttt{+} : lecture et écriture
    \item \texttt{b} : mode binaire
\end{itemize}

Exemples :

\begin{lstlisting}[language=C]
"rb"    // lecture binaire
"w+"    // lecture + écriture
"a+b"   // append + lecture binaire
\end{lstlisting}

\subsection{Tester fopen}

\textbf{Toujours vérifier le retour.}

\begin{lstlisting}[language=C]
FILE *f = fopen("data.txt","r");

if(f == NULL) {
    fprintf(stderr,"Erreur ouverture\n");
} else {
    ...
}
\end{lstlisting}

On peut obtenir des détails avec \texttt{errno} et \texttt{perror()}.

\subsection{Lire et écrire dans un fichier texte}

Fonctions équivalentes à \texttt{printf} et \texttt{scanf} :

\begin{itemize}
    \item \texttt{fprintf}
    \item \texttt{fscanf}
\end{itemize}

\begin{lstlisting}[language=C]
fprintf(f, "%d", x);
fscanf(f, "%d", &x);
\end{lstlisting}

Le fichier est simplement ajouté comme premier argument.

\subsection{Fin de fichier et erreurs}

Fonctions utiles :

\begin{itemize}
  \item \texttt{feof(f)} : fin de fichier (renvoi 0 si pas encore, et un nombre non nulle sinon (uniquement si tentative de lecture échoué en amont))
  \item \texttt{ferror(f)} : erreur de lecture (renvoi 0 si aucune, et un nombre non nulle sinon)
\end{itemize}

Important :

\begin{itemize}
    \item atteindre la fin du fichier n’est pas une erreur
    \item il faut distinguer EOF et erreur
\end{itemize}

\subsection{Fermer un fichier}

\textbf{Toujours fermer un fichier.}

\begin{lstlisting}[language=C]
fclose(f);
\end{lstlisting}

Pourquoi :

\begin{itemize}
    \item les écritures sont bufferisées
    \item \texttt{fclose} garantit que tout est écrit
    \item sinon on peut perdre des données
\end{itemize}

\subsection{printf = fprintf(stdout)}

On peut écrire :

\begin{lstlisting}[language=C]
fprintf(stdout,"Hello");
\end{lstlisting}

équivalent à :

\begin{lstlisting}[language=C]
printf("Hello");
\end{lstlisting}

De même :

\begin{lstlisting}[language=C]
fscanf(stdin,...);
\end{lstlisting}

équivaut à \texttt{scanf}.
(la différence c'est que l'un écrit sur la sortie standard et l'autre sur le fichier spécifié)

\subsection{Choisir dynamiquement la sortie}

On peut stocker le flot dans une variable :

\begin{lstlisting}[language=C]
FILE *out = stdout;
fprintf(out,"Hello");
\end{lstlisting}

Ou bien :

\begin{lstlisting}[language=C]
out = fopen("log.txt","a");
fprintf(out,"Hello");
\end{lstlisting}

Cela permet de choisir dynamiquement la destination.

\subsection{Fichiers texte vs binaires}

\textbf{Texte :}

\begin{itemize}
    \item écrit des caractères
    \item lisible par un humain
    \item plus volumineux
\end{itemize}

\textbf{Binaire :}

\begin{itemize}
    \item écrit directement les octets mémoire
    \item compact
    \item non lisible
\end{itemize}

\subsection{Écriture binaire : fwrite}

Prototype :

\begin{lstlisting}[language=C]
size_t fwrite(const void *ptr,
              size_t taille,
              size_t n,
              FILE *f);
\end{lstlisting}

Exemple :

\begin{lstlisting}[language=C]
int t[12];
fwrite(t, sizeof(int), 12, f);
\end{lstlisting}

Retour :

\begin{itemize}
    \item nombre d’éléments écrits
\end{itemize}

\subsection{Lecture binaire : fread}

\begin{lstlisting}[language=C]
size_t fread(void *ptr,
             size_t taille,
             size_t n,
             FILE *f);
\end{lstlisting}

\begin{lstlisting}[language=C]
int t[12];
fread(t, sizeof(int), 12, f);
\end{lstlisting}

Retour :

\begin{itemize}
    \item nombre d’éléments lus
\end{itemize}

Toujours vérifier ce retour.

\subsection{Position dans un fichier}

\begin{itemize}
    \item \texttt{ftell(f)} : position courante
    \item \texttt{fseek(f, offset, origine)} : déplacer la tête
    \item \texttt{rewind(f)} : revenir au début
\end{itemize}

Origines possibles :

\begin{itemize}
    \item \texttt{SEEK\_SET}
    \item \texttt{SEEK\_CUR}
    \item \texttt{SEEK\_END}
\end{itemize}

\subsection{Flux dans des chaînes}

On peut faire des IO sur des chaînes :

\begin{itemize}
    \item \texttt{sprintf}
    \item \texttt{sscanf}
\end{itemize}

\begin{lstlisting}[language=C]
char buf[100];
sprintf(buf, "%d", x);
\end{lstlisting}

Utile pour :

\begin{itemize}
    \item interfaces graphiques
    \item logs
    \item conversion texte ↔ données
\end{itemize}

\subsection{Bonnes pratiques importantes}

\begin{itemize}
    \item toujours tester \texttt{fopen}
    \item toujours fermer avec \texttt{fclose}
    \item vérifier les tailles lues/écrites
    \item limiter la taille des entrées
\end{itemize}

Sinon :

\begin{itemize}
    \item pertes de données
    \item bugs silencieux
    \item attaques par buffer overflow
\end{itemize}
